\subsection{THETA3-B28: an explicit bound for the actual inverse from source CM evaluations}\label{B28}

The remaining source periods in B26 need not be determined in
order to bound the norm of the actual B17 inverse.  Its finite
CM evaluation matrix gives a lower bound for the source Gram.
For the weight-three source the local constant is explicit,
and the auxiliary source-fibre period cancels.  The statement
is a fixed-input certificate: the CM value matrix, source chart
radii and original measures remain specified inputs.

\paragraph{Data and the exact bound.}
Use B26's source $Y$ of complex dimension $s>0$ and its actual
line $\mathscr L$.  Fix the algebraic coherent basis
$f_1,\ldots,f_r$ of the actual image $E_0$ from B25, and write
\[
 K(g,h)=\sum_{a=1}^r f_a(g)e_a(h),\qquad
 G_{ab}=\langle f_a,f_b\rangle_G,\qquad
 W_{ab}=\int_{[S]}e_a(h)\overline{e_b(h)}\,dh.
 \tag{B28a}
\]
Inner products are conjugate linear in the first variable.
Thus $G$ and $W$ are positive Hermitian matrices, and the
coordinate matrix of the actual $Q_T=TT^*$ is $WG$, as in B17.
No finite restriction operator is substituted for $Q_T$.

Choose distinct source CM points $y_l=[y,b_l]$ that separate
the coefficient sections $e_a$, using the finite Hecke orbit
of the fixed CM point as in B26.  Enlarge the common finite
theta level once for these finitely many points.  Their actual
evaluation matrix has the form
\[
       E_{la}=e_a(1,b_l)={C_x\over t_x}D_{la},
       \qquad D\in M_{L,r}(\overline\Q),\quad
                     \operatorname{rank}D=r.
 \tag{B28b}
\]
The scalar and all finite isogeny factors are those of B25--B26.
In particular $C_x$ and $t_x$ have not been renormalized to one.

At each $y_l$ choose an injective product of oriented disk
charts $|w_j|<\rho_{lj}<1$ whose images in $Y$ are pairwise
disjoint.  After a deeper neat component cover if necessary,
all charts are honest manifold charts.  Its covering degree
is included in the measure below.  Let $\nu_l>0$ be the
conversion from the original adelic source measure to
$\prod_j dx_jdy_j/y_j^2$ on that chart, including its compact
fibre and component weights.  Put
\[
 \kappa_l=\nu_l(4\pi)^s
           \prod_{j=1}^s\left(\rho_{lj}^2-
                                      {\rho_{lj}^4\over2}\right),
 \qquad H=D^*\operatorname{diag}(\kappa_l)D,
 \qquad R=H^{\mathsf t}.
 \tag{B28c}
\]
Both $H$ and $R$ are positive definite.  Then the actual
source Gram and its inverse satisfy the Loewner inequalities
\[
    W\ \geq\ \left|{C_x\over t_x}\right|^2R,
    \qquad
    W^{-1}\ \leq\ \left|{t_x\over C_x}\right|^2R^{-1}.
 \tag{B28d}
\]
For $f=\sum_a q_af_a$, let
$\phi_f=T^*Q_T^{-1}f$.  This is the source vector of least
Hilbert norm among all vectors mapping to $f$.  Its exact
norm and the resulting certificate are
\[
 \|\phi_f\|_S^2=q^*W^{-1}q
       \ \leq\ \left|{t_x\over C_x}\right|^2q^*R^{-1}q.
 \tag{B28e}
\]
Likewise the operator norm on the target with Gram $G$ obeys
\[
 \|Q_T^{-1}\|_{G\to G}\ \leq\
 { |t_x/C_x|^2\over
           \lambda_{\min}(G^{1/2}RG^{1/2})}.
 \tag{B28f}
\]
The inequalities are for the original kernel and its original
measures.  They do not assert that its inverse is algebraic.

\paragraph{Proof of the finite evaluation and the local constant.}
The sections $e_a$ are linearly independent.  Otherwise a
nonzero linear functional on the basis $f_a$ would annihilate
$K(\,\cdot\,,h)$ for every $h$, contradicting the definition
of $E_0$ as the span of these vectors.  B26 proves that finite
translates of $y$ separate the whole finite-dimensional
holomorphic source space, hence a finite list separates the
$e_a$.  Repeated geometric points give proportional evaluation
functionals in a line bundle and may be removed.  The exact
common-scalar formula is B26g; its matrix therefore has rank
$r$.  A finite collection of distinct points on the smooth
quotient admits disjoint sufficiently small product charts.
All their positive radii may be chosen rational.

We compute the chart constant without any statement about
unknown Betti periods.  On one disk first move its centre to
$z=i$ by the indicated real unitary element.  Use the Cayley
coordinate
\[
 z=i{1+w\over1-w},\qquad
 y={1-|w|^2\over|1-w|^2},\qquad
 {dx\,dy\over y^2}={4\,du\,dv\over(1-|w|^2)^2}.
 \tag{B28g}
\]
For a weight-three coefficient write its holomorphic scalar
as $F(z)=(1-w)^3A(w)$.  Any fixed determinant-line comparison
factor is included in $A$; its value at zero is the actual
unitary coefficient at the centre.  The Petersson density is
therefore exactly
\[
       |F(z)|^2y^3{dx\,dy\over y^2}
               =4|A(w)|^2(1-|w|^2)\,du\,dv.
 \tag{B28h}
\]
For several disks the same formula is the product, with
$A$ holomorphic on the polydisk and $A(0)=e(y_l)$ for the
section under consideration.  The central and determinant
characters are constant along these derived-group disks;
their factors are included in this value rather than omitted.

Expand $A(w)=\sum_{\mathbf n\geq0}a_{\mathbf n}w^{\mathbf n}$.
Angular integration kills all mixed monomials.  Every
remaining term in the weighted integral is nonnegative.
The constant term gives
\[
 \int_{|w_j|<\rho_j}|A(w)|^2
                \prod_j(1-|w_j|^2)\,\prod_jdu_jdv_j
 \ \geq\ |A(0)|^2\pi^s
                  \prod_j(\rho_j^2-\rho_j^4/2).
 \tag{B28i}
\]
One may first integrate on smaller closed polydisks where
the power series converges uniformly, and then increase
their radii; positivity and monotone convergence justify
the displayed open-disk formula.  Multiplying by $4^s\nu_l$
and summing the disjoint charts proves
\[
       \|e\|_S^2\ \geq\ \sum_l\kappa_l|e(y_l)|^2.
 \tag{B28j}
\]
This is the weight-three source version of Q11's geometric
evaluation argument, with its sharper radial integral
constant computed explicitly here.

Apply (B28j) to $e=\sum_a u_ae_a$.  Its integral Gram is
$\widehat W=\overline W=W^{\mathsf t}$ in this convention,
so (B28b) gives
$\widehat W\geq|C_x/t_x|^2H$.
Transposition preserves positivity of Hermitian matrices,
and hence gives the first inequality in (B28d).  Conjugating
by $R^{-1/2}$ and inverting its positive eigenvalues gives
the second.  No diagonalization of the unknown source
period matrices was used.

The adjoint formula in B17 gives $Q_T=TT^*$ and its
invertibility on $E_0$.  Therefore $T\phi_f=f$ and
$\phi_f\in\operatorname{im}T^*=(\ker T)^\perp$.
Every other preimage differs by a vector in $\ker T$,
so orthogonality proves the least-norm assertion.  Also
\[
 \|\phi_f\|_S^2
   =\langle Q_T^{-1}f,f\rangle_G
   =q^*G(WG)^{-1}q=q^*W^{-1}q.
\]
This proves (B28e).  Under the isometry given by $G^{1/2}$,
the positive operator $Q_T$ is represented by
$G^{1/2}WG^{1/2}$.  The first inequality of (B28d)
then proves (B28f).\hfill$\square$

\paragraph{A scalar and arithmetic form of the certificate.}
Set $\kappa=\min_l\kappa_l$.  From (B28c),
$\lambda_{\min}(R)\geq\kappa\sigma_{\min}(D)^2$.
Thus the right side of (B28e) is at most
\[
       {|t_x/C_x|^2\over\kappa\sigma_{\min}(D)^2}
                          \sum_a|q_a|^2.
 \tag{B28k}
\]
If $M$ is any nonsingular $r$-row square minor of $D$,
then
\[
 \sigma_{\min}(D)\geq\sigma_{\min}(M)
              \geq { |\det M|\over\|M\|_F^{r-1}}.
 \tag{B28l}
\]
For $r=1$ the denominator is interpreted as one.  The
first assertion follows by adding the other nonnegative
row squares; the second follows by bounding each of the
other $r-1$ singular values by $\|M\|_F$ and using their
product $|\det M|$.

The algebraic determinant estimate of B7 now applies
literally.  For example let $K$ contain the entries of $M$,
$[K:\Q]=D_K$, and let the positive integer $d_0$ make
$d_0M$ integral.  If
$|\sigma(M_{ab})|\leq H_\sigma$ for every other embedding,
the norm of the nonzero algebraic integer
$d_0^r\det M$, together with Hadamard's determinant bound,
gives
\[
 |\det M|\ \geq\ d_0^{-rD_K}
       \prod_{\sigma\ne\mathrm{id}}
                      (r^{r/2}H_\sigma^r)^{-1}.
 \tag{B28m}
\]
This records the required denominator and conjugate bounds;
it does not supply uniform estimates for unspecified input
data.  Source chart radii and the original measure conversion
are likewise not assigned numerical values without an input.

In B26's coherent-source notation the coefficient scalar is
$\gamma=C_x/(t_xt_y)$ and a source CM evaluation brings back
$t_y$.  Formula (B28b) makes their cancellation exact.  The
new bound thus needs no estimate for the extra source fibre
period $t_y$ or the source multiplicity-period matrices.
The retained target theta scalar $C_x/t_x$ can be treated
using the prescribed reference normalization of B22 when
its stated data have been matched.  The algebraic matrices,
integral-frame losses and the exact archive normalization
still have to be supplied.  Nothing here proves the missing
finite-transform derivative rank of A13--A14.

\framework{This is a bound on a specified realization
functional and its genuine adjoint inverse.  The comparison
with algebraic correspondences is the framework of
Grothendieck cote~44, scans~7--8, and cote~16, scans~43--46.
B23--B27 supply the stated algebraic sectors separately;
the present analytic inequality does not itself construct
a new motivic projector.}
